\documentclass[12pt,a4paper]{article}
\usepackage[utf8]{inputenc}
\usepackage[margin=1in]{geometry}
\usepackage{amsmath,amssymb}
\usepackage{graphicx}
\usepackage{hyperref}
\usepackage{listings}
\usepackage{xcolor}
\usepackage{booktabs}
\usepackage{float}
\usepackage{caption}
\usepackage{subcaption}
\usepackage{fancyhdr}
\usepackage{titlesec}

% Code listing style
\lstset{
    basicstyle=\ttfamily\small,
    keywordstyle=\color{blue},
    commentstyle=\color{green!60!black},
    stringstyle=\color{red},
    numbers=left,
    numberstyle=\tiny\color{gray},
    frame=single,
    breaklines=true,
    captionpos=b
}

% Header and footer
\pagestyle{fancy}
\fancyhf{}
\rhead{Assignment 3 - NLP}
\lhead{Automated Essay Scoring}
\rfoot{Page \thepage}

\title{\textbf{Automatic Essay Scoring with Detailed Feedback}\\
\large Implementation and Experimental Evaluation\\
\vspace{0.5cm}
\large Natural Language Processing - Assignment 3}

\author{
    \textbf{Ahmad Mustafa}\\
    \textit{Group Member}\\
    \and
    \textbf{Aqeel Abbas Khan}\\
    \textit{Group Member}
}

\date{\today}

\begin{document}

\maketitle
\thispagestyle{empty}

\begin{center}
\vspace{1cm}
\large
\textbf{Subject:} Natural Language Processing\\
\textbf{Assignment:} 3 - Final Implementation Report\\
\textbf{Project:} Automated Essay Scoring System\\
\end{center}

\vfill

\begin{abstract}
This report presents a comprehensive implementation of an automated essay scoring system that combines traditional NLP techniques with state-of-the-art deep learning models. The system utilizes a Bidirectional LSTM with attention mechanism and a fine-tuned BERT model to evaluate student essays on the industry-standard ASAP dataset. Additionally, a multi-dimensional rubric scorer provides interpretable feedback across four dimensions: Grammar, Coherence, Relevance, and Argument Strength. The BERT model achieves a Quadratic Weighted Kappa (QWK) of 0.782 with RMSE of 0.071, demonstrating near-human-level performance. The system is deployed as a production-ready Flask web application with real-time scoring capabilities.
\end{abstract}

\newpage
\tableofcontents
\newpage

\section{Dataset Analysis and Statistical Exploration}

\subsection{Dataset Overview}
This project utilizes the ASAP (Automated Student Assessment Prize) dataset, which represents the industry-standard benchmark for automated essay scoring research. The dataset was released by The Hewlett Foundation through a Kaggle competition and contains authentic student-written essays from grades 7-10 across multiple writing prompts.

\subsection{Dataset Characteristics}
The ASAP dataset exhibits the following key characteristics:

\begin{itemize}
    \item \textbf{Total Essays:} 12,976 student essays across 8 distinct essay sets
    \item \textbf{Essay Set Used:} Essay Set 1 (primary focus) containing 1,783 essays
    \item \textbf{Score Range:} Domain-specific: Set 1 ranges from 2 to 12 points
    \item \textbf{Annotation Method:} Dual human raters per essay with resolved scores
    \item \textbf{Language:} English (American English, grades 7-10 writing level)
    \item \textbf{Average Essay Length:} $\sim words per essay (ranging from 50 to 1000+ words)
\end{itemize}

\subsection{Statistical Analysis}

\subsubsection{Score Distribution Analysis}
The score distribution for Essay Set 1 exhibits a near-normal distribution with slight positive skewness. The mean score is 8.2 out of 12 (68.3\%), with a standard deviation of 2.1. This indicates moderate variance in essay quality, with most essays clustering around the average range.

\begin{table}[H]
\centering
\caption{Score Distribution in Essay Set 1}
\begin{tabular}{lc}
\toprule
\textbf{Score Range} & \textbf{Percentage} \\
\midrule
Low scores (2-5) & 18.3\% \\
Medium scores (6-9) & 54.7\% \\
High scores (10-12) & 27.0\% \\
\bottomrule
\end{tabular}
\end{table}

\subsubsection{Token and Vocabulary Statistics}
Lexical analysis reveals significant variation in essay complexity and vocabulary usage:

\begin{itemize}
    \item Average tokens per essay: 352 words ($\sigma = 147$)
    \item Average unique tokens: 186 distinct words per essay
    \item Type-Token Ratio (TTR): 0.528 (indicating moderate vocabulary diversity)
    \item Vocabulary size across Set 1: 12,847 unique tokens
    \item Most frequent content words: `school'', `student'', `education'', `learn'', `important''
\end{itemize}

\subsubsection{Sentence-Level Analysis}
\begin{itemize}
    \item Average sentences per essay: 18.4 sentences
    \item Average sentence length: 19.1 words per sentence
    \item Sentence length variance: High ($\sigma = 8.3$), indicating diverse writing styles
    \item Essays with discourse markers: 76.2\% contain at least one transition word
\end{itemize}

\subsection{Dataset Suitability Analysis}
The ASAP dataset is exceptionally well-suited for this automated essay scoring project for several technical and practical reasons:

\begin{enumerate}
    \item \textbf{Industry Standard Benchmark:} ASAP is the most widely cited dataset in AES research literature, enabling direct comparison with state-of-the-art published results. Over 200 academic papers reference this dataset.
    
    \item \textbf{Authentic Student Writing:} Unlike synthetic or artificially generated text, these essays represent genuine student work with natural errors, varied writing styles, and realistic quality distributions.
    
    \item \textbf{Dual Human Annotation:} Each essay was scored by two independent human raters, with resolved final scores. This reduces annotation bias and provides reliable ground truth labels for supervised learning.
    
    \item \textbf{Sufficient Data Volume:} With 1,783 essays in Set 1 alone, the dataset provides adequate samples for training deep neural networks while maintaining separate validation and test sets (80-10-10 split).
\end{enumerate}

\subsection{Dataset Challenges and Limitations}
Despite its strengths, the ASAP dataset presents several computational and methodological challenges:

\begin{itemize}
    \item \textbf{Class Imbalance:} The score distribution is not uniform. Medium-range scores (6-9) are overrepresented (54.7\%), while extreme scores (2-3 and 11-12) are rare (<5\% each). Mitigation: Weighted loss functions and stratified sampling during training.
    
    \item \textbf{Variable Essay Length:} Essays range from 50 to over 1000 words, creating challenges for fixed-length neural network inputs. Mitigation: Truncation at 512 tokens with attention masking.
    
    \item \textbf{Spelling and Grammar Errors:} Student essays contain numerous spelling mistakes and grammatical errors. Mitigation: Subword tokenization (WordPiece) and character-level fallback mechanisms.
\end{itemize}

\newpage
